\section{Related Work}

\subsection{Alignment and Preference Optimization}
Optimizing Large Language Models (LLMs) to align with human intent has predominantly relied on Reinforcement Learning from Human Feedback (RLHF) \cite{ouyang2022training}, utilizing Proximal Policy Optimization (PPO) \cite{schulman2017proximal}. Recently, Direct Preference Optimization (DPO) \cite{rafailov2023direct} emerged as a stable, offline alternative that mathematically bypasses the need for an explicit reward model by optimizing the policy directly over preference pairs. However, both PPO and DPO traditionally rely on human-annotated preference data, which is expensive to scale and prone to subjective stylistic biases.

\subsection{Hallucination Mitigation in LLMs}
In knowledge-intensive tasks, such as educational explanation generation, factual correctness is paramount. Prior work has explored grounding LLMs using retrieval-augmented generation (RAG) \cite{lewis2020retrieval} or knowledge graphs. In the educational domain, techniques have been developed to enhance logical reasoning \cite{bao2023enhancing, bao2024abstract} and automate the generation of scientific claims from multiple-choice questions \cite{bao2023multi2claim}. Furthermore, there has been significant progress in predicting and analyzing the quality of learnersourced multiple-choice questions \cite{bao2022deepqr}, and employing Large Language Models for iterative enhancement of generated explanations \cite{bao2025exploring}. Recent studies have also demonstrated the utility of Natural Language Inference (NLI) models \cite{honovich2022true} to detect and penalize hallucinations, or even training Transformers to explain unexpected inputs \cite{bao2022abductionrules}. \cite{gao2023enabling} demonstrated the utility of NLI for self-correction, but purely logical optimization often degrades the fluency and readability of the text---a phenomenon we define as the "alignment tax."

\subsection{LLM-as-a-Judge and Evaluation Bias}
The advent of powerful instruction-tuned models like GPT-4 has popularized their use as automated evaluators \cite{zheng2024judging}. However, recent studies highlight significant biases in LLM-as-a-judge paradigms, including positional bias, self-preference bias, and verbosity bias \cite{wang2023large}. Our research corroborates these findings, demonstrating that leading commercial models systematically penalize concise, factually dense explanations in favor of verbose, hallucinated rhetoric. Our Hybrid-DPO framework circumvents this by relying on objective, deterministic metrics (NLI) fused with constrained fluency verifiers during the alignment phase.
